\section{三类问题的求解结果与成本}
\label{sec:opt-mmb-results-costs}

前三节分别处理了三类问题，用到的工具有很多重叠：梯度、投影、乘子、二次上界。工具相同容易让人以为它们的结果也可以用同一把尺子衡量，例如把“共同下降方向的范数小于$10^{-6}$”“外梯度残差小于$10^{-6}$”“超梯度小于$10^{-6}$”都称为“已收敛”。本节以“得到并检验一个规定的答案”为共同任务，先核对三类输出各自满足什么关系，再核对为得到它付出了什么。三类问题的残差数值含义不同，不作相互排序。

\subsection{不同最优性标准}
\label{sec:opt-mmb-standards}

\paragraph{各自改变了标量优化的哪一项}
标量优化由三部分组成：候选集合、比较候选的规则、以及计算目标值的方式。多目标问题保留了候选集合和目标的计算方式，改变的是比较规则：向量之间只有支配这一偏序，答案因而是一个集合，选出其中一点需要额外的偏好。极小极大问题保留了数值比较，改变的是目标值的计算方式：$\bar\phi(\bm x)$本身是一个优化问题的最优值，而且依赖于双方的选择次序。双层问题同样改变目标值的计算方式，但内层有自己的目标，外层值依赖于内层解如何选出；选解规则不同，外层目标就是不同的函数。三类问题都会用到梯度和乘子，这并不让它们共享一个解概念。

\paragraph{局部判据、全局证书与内层精度}
表~\ref{tab:opt-mmb-standards}逐项列出三类问题的局部判据能说明什么、有没有可以计算的全局证书，以及判据对内层计算的精度要求。

\begin{table*}[!htbp]
  \centering
  \small
  \begin{tabularx}{\linewidth}{@{}p{12mm}XXXX@{}}
    \toprule
    问题 & 局部判据 & 判据能支持的结论 & 可计算的全局证书 & 对内层精度的要求 \\
    \midrule
    多目标 & 梯度凸包的最小范数$\lVert\bar{\bm g}\rVert_2$ & 不存在一阶共同下降方向；各指标凸时为某个加权和的最小点，正权时为Pareto最优 & 凸情形下加权和的最优性证书；有限候选可逐一比较；阈值问题可用第~\ref{chap:constrained-optimization}章的对偶界 & 无内层；回溯须对每个指标检查 \\
    极小极大 & 投影残差$\lVert\bm u-\Pi_K(\bm u-\eta\bm F(\bm u))\rVert_2/\eta$ & 凸凹且可微时残差为零等价于鞍点；非凸时只说明是驻点 & 间隙$\bar\phi(\bm x)-\underline\phi(\bm y)$，须求两侧最优响应或其界 & 内层最大化不精确时只得到$\bar\phi$的下界 \\
    双层 & 超梯度范数$\lVert\nabla J_T\rVert_2$或$\lVert\nabla\widehat J\rVert_2$ & 对所选目标（有限步或局部分支）的一阶驻点 & 一般没有；还须另行核验选解规则与分支 & 内层残差与伴随残差共同决定超梯度误差，并被小曲率放大 \\
    \bottomrule
  \end{tabularx}
  \caption{三类问题的最优性标准。各行的判据针对不同对象，同样小的数值代表不同强度的结论。}
  \label{tab:opt-mmb-standards}
\end{table*}

表中第二、三列的差距是本章反复出现的主题。共同下降法停在$\bar{\bm g}=\bm 0$，可能是一个被支配的点（第~\ref{sec:opt-mmb-common-descent}小节的$f_1=f_2=-x^2$）；极小极大的零梯度点可能有正的间隙（$\phi=y^2-x^2$在原点的间隙为$2$）；双层问题的超梯度为零，可能只是局部分支上的驻点，而真正的选解在那里跳变（倾斜双势阱在$\theta=0$处）。因此报告结果时，应写明“满足了哪一个判据、对哪一个目标、在什么条件下它意味着什么”，而不是只写一个“已收敛”。

\subsection{同一输出要求下的计算代价}
\label{sec:opt-mmb-costs}

\paragraph{一轮外层迭代包含什么}
三类方法的一轮迭代，实际工作量差别很大。表~\ref{tab:opt-mmb-costs}以梯度、Hessian--向量积和投影为单位记录每轮代价，其中$m$是指标个数，$T$是内层步数，$k$是伴随方程的共轭梯度迭代次数。

\begin{table*}[!htbp]
  \centering
  \small
  \begin{tabularx}{\linewidth}{@{}p{30mm}XX@{}}
    \toprule
    方法 & 每轮主要计算 & 额外存储或开销 \\
    \midrule
    加权和、阈值约束 & 一个标量或约束优化问题的完整求解（每个权重或阈值一次） & 描出前沿需对许多权重或阈值重复求解 \\
    共同下降法 & $m$个梯度；$m$维单纯形二次规划（形成Gram矩阵需$O(m^2d)$）；每次回溯试探求$m$个指标值 & 存储$m$个梯度 \\
    同时梯度下降—上升 & 一组偏梯度，$X$、$Y$上各一次投影 & 报告间隙需另解两个最优响应问题 \\
    外梯度法 & 两组偏梯度，$X$、$Y$上各两次投影 & 同上；平均保证需累计试探点 \\
    迭代展开（反向） & $T$步内层更新；反向$T$次Hessian--向量积与混合导数乘积 & 存储$O(Td)$，或用检查点以重算换存储 \\
    迭代展开（前向） & $T$步内层更新，每步$p$次Hessian--向量积 & 存储$O(dp)$ \\
    隐式微分 & 内层求解到容差；约$k$次Hessian--向量积解伴随方程；一次混合导数乘积 & 存储$O(d)$；须检查Hessian正定 \\
    \bottomrule
  \end{tabularx}
  \caption{各方法一轮迭代的主要代价。双层方法的每一轮外层更新都包含一次完整的内层求解或展开。}
  \label{tab:opt-mmb-costs}
\end{table*}

“一轮外层更新”在这些方法中代表的工作量并不相同。外梯度一轮相当于两轮同时梯度下降—上升；双层方法的一轮外层更新，包含了一次内层优化（$T$步或求解到容差）再加上一次反向或伴随计算，可能是外层梯度步本身代价的成百上千倍。按轮数比较不同方法没有意义。

\paragraph{先固定输出要求，再比较路线}
比较求解路线之前，先固定要返回的东西及其容差：是只要可行性，还是要某个局部残差小于$\varepsilon$，还是要一个全局证书（例如间隙不超过$\varepsilon$）。然后在同一问题、同一要求下比较总代价。近似的内层计算必须按它对最终输出的影响计费：若内层只解到使超梯度误差为$\delta$的精度，那么由引理~\ref{lem:opt-mmb-inexact-outer}，外层最多只能把驻点残差压到$\delta$量级，再多的外层迭代都是浪费；反之，为了一个只需粗略方向的早期外层步把内层解到机器精度，同样是浪费。极小极大问题中，若最终要报告间隙，求两侧最优响应的代价可能与求解本身相当，也必须算入。

本章的三类问题都是静态的：比较、响应和嵌套都发生在一次决策之内。若决策随时间展开，后一步的选择依赖于前一步观察到的结果；若对象在反馈下运行，控制量又会改变后续状态；若多方同时行动且各有信息，最优响应要换成均衡概念。这些结构在卷五中依次加入，那里的博弈部分会在本章极小极大与最优响应的基础上研究多方响应及其演化，可参考\BookChapterRef{chap:game-theory-and-multi-agent-decision}{卷五“多方策略的响应与演化”章}。
